\subsection{Full local fixed groups from primitive valuation uniqueness}

The following deduction uses the uniqueness assertion of
\citet[Theorem A; Proposition 2.4(vii)]{D_MochizukiTsujimura2024}.
It is a consequence of that opened primary-source theorem, with the valuation
and stabilizer arguments supplied below. It does not assert a new
resolution-of-nonsingularities theorem.

\begin{claim}{C.7h}{Proved}
Let $K/\Q_p$ be finite, let $U/K$ be a smooth geometrically connected
hyperbolic curve, and fix an algebraic closure $\bar K$. Write
\[
 1\longrightarrow\Pi\longrightarrow E_U\longrightarrow G_K
 \longrightarrow1,\qquad
 H=G_{K(\mu_\infty)}\subset G_K,
\]
where $E_U=\pi_1(U)$ and $\Pi=\pi_1(U_{\bar K})$ are the full
profinite groups. Choose compatible basepoints and paths, so that the
sections and subgroups below are actual subgroups rather than only
conjugacy classes. For every open subgroup $H'\subset H$ the following
hold.

If $x\in U(K)$ and $s_x:G_K\to E_U$ is its continuous point section,
then
\[
 \Pi^{s_x(H')}=C_\Pi(s_x(H'))=1.
\]
If $c\in\overline U\setminus U$ is $K$-rational, $a$ is any nonzero
$K$-rational tangent vector at $c$, and $s_{c,a}:G_K\to E_U$ is the
associated continuous tangential section, then
\[
 \Pi^{s_{c,a}(H')}=C_\Pi(s_{c,a}(H'))=I_c
 \simeq\widehat\Z(1),
\]
where $I_c$ is the full geometric cusp inertia subgroup determined by the
chosen tangential basepoint. More generally the same equality holds for
every continuous cuspidal section $s:G_K\to D_c\subset E_U$, where
$D_c$ is the arithmetic decomposition group of a chosen compatible lift
of the rational cusp $c$; its Kummer cocycle need not come from a tangent
coefficient in $K^*$. Thus the same equalities hold for elements
fixed by some open subgroup of $H$. They hold for every residue prime,
including $2$ and $3$, with no reduction hypothesis.

In particular, on
$\mathbf P^1_{\Q_5}\setminus\{0,1,\infty\}$ the full fixed group of
the point $2$ is trivial, and the full fixed group of the tangent at $0$
with coefficient $1$ is $I_0$. Together with C.6, this determines both
halves of C4.
\end{claim}
\begin{proof}
\emph{The exact uniqueness theorem and its field.}
Take $\Sigma$ to be the set of all primes. Let
$\widetilde U\to U$ be the universal geometrically pro-$\Sigma$
covering, so that its group of automorphisms over $U$ is $E_U$.
Its function field
\[
 L=K(\widetilde U)
\]
is the algebraic union of the function fields of its connected finite
etale subcoverings. It contains $\bar K$, and
$\Pi=\operatorname{Gal}(L/\bar K(U))$. This is an algebraic function
field extension, without any completion. The completed algebraic closure
$\Omega$ that occurs elsewhere in the cited paper is used for analytic
point evaluation; the valuation bijection and uniqueness assertion used
here are stated on $K(\widetilde U)$ itself.

In Definition 2.2(i) of the cited paper, a $p$-valuation on $L$ means a
valuation ring $R\subset L$ satisfying $R\cap K=\mathcal O_K$.
It is primitive if its maximal ideal is the only prime ideal of $R$
containing $p$. No rank-one hypothesis is imposed. Proposition 2.3(v)
identifies all such valuations, and their primitive subcollection, with
the space $\operatorname{VE}(\widetilde U)$ of center chains and its
primitive subspace. The bijection is equivariant for $E_U$: an
automorphism carries each valuation center on a compactified model to
the center of the transported valuation on the transported model.
In particular the bijection includes valuations centered at cusps of
compactifications, even though those are not called point-theoretic
points of the open curve in Definition 2.2(ii).

Theorem A applies to every $U/K$ under consideration and gives
$\Sigma$-RNS. Proposition 2.4(vii) then states that, for a closed subgroup
$J\subset G_K$ such that $J\cap I_K$ admits a continuous surjection
onto $\Z_\ell$ for some $\ell\ne p$, the image of any continuous
section $s:J\to E_U$ fixes at most one primitive valuation on $L$.
The additional open-image hypothesis on the $\ell$-adic cyclotomic
character is used only for existence, and is not a hypothesis of this
uniqueness assertion. Nor is $J$ required to be normal. We will supply a
fixed primitive valuation directly, so only uniqueness is needed.
The proof of the cited proposition makes this distinction explicit:
existence uses its Proposition 2.3(xi), whereas uniqueness uses the
bounded-distance conclusion of
\citet[Corollary 1.15(iv)]{BC_HMTopicsIV} together with primitive
center-chain separation. The latter corollary requires an open tame
inertia image; its extra cyclotomic-fullness condition belongs to
part (iii), concerning existence, and is not a hypothesis of part (iv).

\emph{The subgroup hypothesis.}
Let $q$ be the residue cardinality of $K$, let $P_K$ be wild inertia,
and write
\[
 T=I_K/P_K\simeq\prod_{\ell\ne p}\Z_\ell(1).
\]
The tame quotient $G_K/P_K$ is $T\rtimes\widehat\Z$, with an
arithmetic Frobenius acting on $T$ by multiplication by $q$. Hence the
image of the closed commutator subgroup of $G_K$ in $T$ contains
$(q-1)T$: lift a tame element $t$ and a Frobenius element, and compute
their commutator in this quotient. Every cyclotomic character has
abelian target, so $[G_K,G_K]\subset H$. Also $H\subset I_K$:
the prime-to-$p$ roots of unity generate the maximal unramified extension
of $K$, since every finite residue extension is generated by its
Teichmuller roots of unity. Consequently the image of $H$ in each
$\Z_\ell(1)$, $\ell\ne p$, contains
$(q-1)\Z_\ell(1)$ and is open.

The subgroup $H'$ is closed in $G_K$, because it is an open subgroup of
the compact group $H$. Its image in $\Z_\ell(1)$ has finite index in
the image of $H$: the index of an image is at most $[H:H']$.
It is therefore an open nonzero subgroup of $\Z_\ell(1)$, hence
isomorphic as a topological group to $\Z_\ell$. The map onto its image,
followed by this isomorphism, is the required continuous surjection
\[
 H'\cap I_K=H'\twoheadrightarrow\Z_\ell.
\]
An identification of that image with $\Z_\ell$ rescales a lattice; no
surjectivity onto a preselected tame generator is required by the
theorem. This argument applies to arbitrary open $H'$, without requiring
$H'$ to be normal.

\emph{A primitive composite valuation at the chosen lift.}
Write $b=x$ in the ordinary case and $b=c$ in the cuspidal case.
The chosen basepoint determines a compatible lift $\widetilde b$ of
$b$ to the smooth compactifications of all the finite subcoverings of
$\widetilde U$. The point section or cuspidal section preserves this
lift. Let $A\subset L$ be its horizontal valuation ring, obtained by
taking the union of the local rings at these compatible points, and let
$\mathfrak m_A$ be its maximal ideal. At each finite geometric level
the point has algebraic residue; the constants of $L$ contain $\bar K$.
Thus the residue field of $A$ is precisely $\bar K$, and evaluation gives
\[
 A\twoheadrightarrow\bar K,
 \qquad\ker(\operatorname{ev}_{\widetilde b})=\mathfrak m_A.
\]
This construction is equally valid at a boundary point, allowing the
ramification indices of the compactified covers there. Define
\[
 R=\{f\in A:\operatorname{ev}_{\widetilde b}(f)
                     \in\mathcal O_{\bar K}\}.
\]
It is a valuation ring: it is the composition of the horizontal valuation
with the unique extension of the $p$-adic valuation to its residue field
$\bar K$. Explicitly, a horizontal pole has inverse in
$\mathfrak m_A\subset R$, and a horizontal unit with negative
$p$-adic residue valuation has inverse in $R$. Moreover
$R\cap K=\mathcal O_K$ and
\[
 R/\mathfrak m_A=\mathcal O_{\bar K}.
\]
The element $p$ is a unit in $A$, so
$\mathfrak m_A=p\mathfrak m_A\subset pR$. Every prime ideal of $R$
containing $p$ therefore contains $\mathfrak m_A$. In
$\mathcal O_{\bar K}$ one has
$\sqrt{p\mathcal O_{\bar K}}=\mathfrak m_{\bar K}$: a positive
valuation becomes at least $v(p)$ after multiplication by a sufficiently
large positive integer. It follows that the only prime ideal of $R$
containing $p$ is its maximal ideal. Thus $R$ is primitive in exactly
the sense required by Proposition 2.4(vii), even though this composite
valuation has a nontrivial horizontal rank.

Crucially, the horizontal place is recovered from this primitive
valuation by
\[
 R[1/p]=A.                                                   \tag{C.7h.1}
\]
Indeed $p$ is a horizontal unit, so the left side is contained in $A$.
For $f\in A$ with nonzero residue, choose $n\ge0$ such that
$p^n\operatorname{ev}_{\widetilde b}(f)$ is integral; then $p^nf\in R$.
If the residue is zero, $f$ already belongs to $R$.
Thus (C.7h.1) remembers the actual horizontal point, not merely its
reduction on one model.

The whole point or cuspidal section $s:G_K\to E_U$ fixes $R$.
It preserves the horizontal lift, and on the residue field acts through
the given element of $G_K$, which preserves the unique extended
$p$-adic valuation of $\bar K$.

\emph{Centralizers and horizontal stabilizers.}
Let $g\in C_\Pi(s(H'))$. Since $R$ is fixed by $s(H')$, its transport
$gR$ is also fixed: for every $\sigma\in H'$,
\[
 s(\sigma)(gR)=g(s(\sigma)R)=gR.
\]
The valuation $gR$ is again primitive, since $g$ fixes $K$ and permutes
prime ideals. The uniqueness theorem gives $gR=R$. As $g(p)=p$,
(C.7h.1) gives $gA=A$, so $g$ stabilizes the compatible horizontal
lift $\widetilde b$.

At an ordinary point the stabilizer of $\widetilde x$ in $\Pi$ is
trivial. Indeed, on every finite connected geometric Galois cover the
cover is etale above $x$, and its residue fields there are $\bar K$.
Its decomposition group is therefore trivial. Such covers form a
cofinal system of finite quotients of $\Pi$, so an element stabilizing
all these lifts is the identity.

At a cusp the stabilizer of $\widetilde c$ in $\Pi$ is precisely its
geometric inertia $I_c$: the geometric residue field is algebraically
closed, so decomposition and inertia coincide. The characteristic-zero
local parameter covers identify this full cusp inertia with
$\widehat\Z(1)$; its injection into the geometric hyperbolic-curve
group is the cusp injection recalled in B.3
\citep[Proposition 1.2]{BC_MochComb}. We have proved
$C_\Pi(s(H'))\subset I_c$ for every section into this $D_c$.
Conversely, conjugation by every element of $D_c$ on $I_c$ is the
cyclotomic action: locally it sends a root-of-unity multiplier in
$z^{1/n}$ to its Galois conjugate; the inertia component of a section
cocycle acts trivially because $I_c$ is abelian.
Since every element of $H'$ fixes
all roots of unity, this action is trivial. Hence
$I_c\subset C_\Pi(s(H'))$, proving equality.

Every comparison used here concerns the full arithmetic extension and
its actual restricted sections. In particular no lift of a fixed
element through a prime-to-$p$ quotient, and no detection assertion for
nilpotent or genus-zero quotients, is used.
\end{proof}

The proof also gives uniqueness of the underlying ordinary or cuspidal
point when two geometric point or tangential sections become conjugate
over a finite extension of $K(\mu_\infty)$: apply the same primitive
valuation uniqueness to the common restricted section, after realizing
the finite data over a finite extension of $K$. This does not assert
uniqueness of the tangent vector at a fixed cusp.
